\documentclass[11pt]{article}

\usepackage[T1]{fontenc}
\usepackage{lmodern}
\usepackage{amsmath,amssymb,amsthm,mathtools}
\usepackage{enumitem}
\usepackage[margin=1.12in]{geometry}
\usepackage{microtype}
\usepackage[hidelinks]{hyperref}

\newtheorem{theorem}{Theorem}[section]
\newtheorem{lemma}[theorem]{Lemma}
\newtheorem{proposition}[theorem]{Proposition}
\newtheorem{corollary}[theorem]{Corollary}
\theoremstyle{definition}
\newtheorem{definition}[theorem]{Definition}

\newcommand{\NN}{\mathbb N}
\newcommand{\PP}{\mathbb P}
\newcommand{\End}{\operatorname{End}}
\newcommand{\rank}{\operatorname{rank}}
\newcommand{\Mat}{\operatorname{M}}
\newcommand{\Free}{k\langle X_0,X_1,X_2\rangle_+}
\newcommand{\coeff}[2]{[s^{#1}]#2}

\title{A Counterexample to K\"othe's Conjecture}
\author{}
\date{}

\begin{document}
\maketitle

\begin{abstract}
For every countable algebraically closed field $k$, there exists a three-generated nil nonunital $k$-algebra $A$ such that $\Mat_2(A)$ is not nil.  Consequently, there exists a unital ring containing two nil left ideals whose sum is not nil.
\end{abstract}

\section{Introduction}\label{sec:introduction}

K\"othe's 1930 question asks whether an associative ring with no nonzero nil two-sided ideal can contain a nonzero nil one-sided ideal \cite{Kothe1930}.  Equivalently, the sum of two nil left ideals of an arbitrary ring must be nil.  In 1956 Amitsur proved that, for every ring $R$, the Jacobson radical of $R[x]$ has the form $I[x]$ for a nil ideal $I\subseteq R$ \cite{AmitsurRadicals1956}.  He also proved that if $A$ is a nil algebra over an uncountable field, then $A[x_1,\ldots,x_n]$ is nil for every finite set of commuting indeterminates \cite{AmitsurInfinite1956}.

Krempa proved in 1972 that K\"othe's conjecture is equivalent, among other statements, to each of the following: $M_2(N)$ is nil for every nil ring $N$, and $N[x]$ is Jacobson radical for every nil ring $N$ \cite{Krempa1972}.  The matrix formulation was obtained independently by Sands \cite{Sands1973}.  Smoktunowicz later constructed a nil algebra over a countable field whose polynomial ring is not nil \cite{Smoktunowicz2000}, thereby refuting the stronger assertion that $N[x]$ must be nil for every nil ring $N$; this did not settle K\"othe's conjecture.  Beidar, Fong, and Puczy\l{}owski proved that $N[x]$, for $N$ nil, cannot be homomorphically mapped onto a ring containing a nonzero idempotent \cite{BeidarFongPuczylowski2001}.  Smoktunowicz subsequently proved that every primitive ideal of $N[x]$ is of the form $I[x]$ for an ideal $I$ of $N$, and deduced the equivalent formulation that K\"othe's conjecture holds if and only if polynomial rings over nil rings are never right primitive \cite{Smoktunowicz2005}.

In September 2026 a machine-checked Lean~4 proof established a counterexample to Krempa's matrix formulation.  It was generated by a pre-release GPT-6 Astra model in an Epoch AI LeanOpenProblems evaluation and is archived, together with verification data, in the repository maintained by Adamczewski \cite{Adamczewski2026}.  The construction proved below is derived from that formal proof and is not an independent discovery of the counterexample.  The presentation reorganizes the argument in ordinary mathematical language and replaces some formal ingredients by standard arguments; in particular, the common-zero step is deduced from a Segre--Veronese embedding and the projective dimension theorem, and the finite linearization is written explicitly.

\section{Statement and notation}\label{sec:statement}

An element $x$ of a ring is \emph{nilpotent} if $x^n=0$ for some $n\ge1$.  A ring, algebra, or one-sided or two-sided ideal is \emph{nil} if each of its elements is nilpotent.

All rings are associative.  A \emph{nonunital $k$-algebra} is a $k$-vector space equipped with an associative $k$-bilinear multiplication, without a specified identity.  If $B$ is a ring, $B[s]$ denotes the polynomial ring in a central indeterminate $s$.

\begin{theorem}[Main theorem]\label{thm:main}
Let $k$ be a countable algebraically closed field.  There exists a three-generated nil nonunital $k$-algebra $A$ such that $\Mat_2(A)$ contains a nonnilpotent element.  Consequently there exists a unital ring containing two nil left ideals whose sum is not nil.
\end{theorem}

\section{A projective common-zero lemma}\label{sec:projective}

Let $N\ge1$, and for each $1\le j\le N$ let
\[
y^{(j)}=(y^{(j)}_0,y^{(j)}_1,y^{(j)}_2)
\]
be a block of three variables.  A polynomial in all these variables is \emph{multihomogeneous of degree $r$ in every block} if each monomial occurring with nonzero coefficient has total degree $r$ in the variables belonging to each fixed block.  The zero polynomial is regarded as multihomogeneous of every multidegree.

\begin{lemma}[Projective common-zero lemma]\label{lem:projective}
Let $k$ be algebraically closed, let $N,r\ge1$, and let
\[
f_1,\ldots,f_q\in k[y_i^{(j)}:1\le j\le N,\ 0\le i\le2]
\]
be multihomogeneous of degree $r$ in every block.  If $q\le2N$, then there exist
\[
v^{(1)},\ldots,v^{(N)}\in k^3\setminus\{0\}
\]
such that
\[
f_a(v^{(1)},\ldots,v^{(N)})=0
\qquad(1\le a\le q).
\]
\end{lemma}

\begin{proof}
Set
\[
X=(\PP^2_k)^N,
\qquad \dim X=2N.
\]
Let $\mathcal M_r$ be the set of monomials of total degree $r$ in three variables, so
\[
|\mathcal M_r|=\binom{r+2}{2}.
\]
The product of the degree-$r$ Veronese embeddings followed by the Segre embedding gives a closed immersion
\[
\nu:X\longrightarrow\PP^{M-1}_k,
\qquad
M=|\mathcal M_r|^N,
\]
whose homogeneous coordinates are indexed by tuples $(m_1,\ldots,m_N)\in\mathcal M_r^N$ and are
\[
\prod_{j=1}^N m_j(y^{(j)}).
\]
Thus every polynomial of multidegree $(r,\ldots,r)$ is the pullback under $\nu$ of a homogeneous linear form on $\PP^{M-1}$.

For each $a$ let $\ell_a$ be the corresponding linear form, and let
\[
L=V(\ell_1,\ldots,\ell_q)\subseteq\PP^{M-1}.
\]
The common kernel of $q$ linear forms in an $M$-dimensional vector space has dimension at least $M-q$.  Since
\[
M\ge3^N>2N\ge q,
\]
this kernel is nonzero, and therefore $L$ is a nonempty projective linear subspace satisfying
\[
\dim L\ge M-1-q.
\]
Consequently
\[
\dim\nu(X)+\dim L
\ge 2N+M-1-q
\ge M-1.
\]
Both $\nu(X)$ and $L$ are projective varieties.  By the projective dimension theorem, $\nu(X)\cap L\ne\varnothing$; see \cite[Chapter I, Theorem 7.2]{Hartshorne}.  A point of $X$ is represented by $N$ nonzero triples in $k^3$, and belonging to $L$ says precisely that all the $f_a$ vanish.  This proves the lemma.
\end{proof}

\section{Finite linearization of a noncommutative polynomial}\label{sec:linearization}

Let
\[
\Free=k\langle X_0,X_1,X_2\rangle_+
\]
denote the free nonunital associative $k$-algebra on three generators.  Thus its $k$-basis is the set of all nonempty words in $X_0,X_1,X_2$.

\begin{lemma}[Finite linearization]\label{lem:linearization}
Let $\widetilde B$ be a unital $k$-algebra, let $B\subseteq\widetilde B$ be a nonunital $k$-subalgebra, and let $a_0,a_1,a_2\in B$.  Let $f\in\Free$.  Then there exist an integer $d\ge1$ and matrices
\[
C_0,C_1,C_2,D_0,D_1,D_2\in\Mat_d(k)
\]
such that each $D_i$ has all rows zero except possibly its first row, with the following property.

Put
\begin{equation}\label{eq:linearized-operator}
T_f(a;s)=\sum_{i=0}^2(C_i+sD_i)a_i\in\Mat_d(\widetilde B[s]),
\end{equation}
where scalar matrices act by scalar multiplication on the entries $a_i$.  If $T_f(a;s)$ is nilpotent, then $f(a_0,a_1,a_2)$ is nilpotent in $B$.
\end{lemma}

\begin{proof}
If $f=0$, take $d=1$ and $C_i=D_i=0$ for $i=0,1,2$.  We may therefore assume $f\ne0$.  Write $f$ as a finite sum of distinct nonempty words,
\begin{equation}\label{eq:fwords}
f=\sum_{\alpha=1}^m c_\alpha
X_{i_{\alpha,1}}X_{i_{\alpha,2}}\cdots X_{i_{\alpha,\ell_\alpha}},
\end{equation}
where $c_\alpha\in k\setminus\{0\}$ and $\ell_\alpha\ge1$.

Introduce one distinguished index, denoted $0$.  For every word with $\ell_\alpha\ge2$, introduce additional indices
\[
(\alpha,2),(\alpha,3),\ldots,(\alpha,\ell_\alpha).
\]
Let $\Omega$ be the resulting finite index set; choose any ordering with the distinguished index first, and put $d=|\Omega|$.

Define a matrix $T(s)\in\Mat_d(\widetilde B[s])$ by the following nonzero entries:
\begin{align}
T_{0,0}(s)
&=s\sum_{\ell_\alpha=1}c_\alpha a_{i_{\alpha,1}},\label{eq:T00}\\
T_{0,(\alpha,2)}(s)
&=s c_\alpha a_{i_{\alpha,1}}
\qquad(\ell_\alpha\ge2),\label{eq:Troot}\\
T_{(\alpha,j),(\alpha,j+1)}(s)
&=a_{i_{\alpha,j}}
\qquad(2\le j<\ell_\alpha),\label{eq:Tchain}\\
T_{(\alpha,\ell_\alpha),0}(s)
&=a_{i_{\alpha,\ell_\alpha}}
\qquad(\ell_\alpha\ge2).\label{eq:Tlast}
\end{align}
All other entries are zero.  Every entry is a $k$-linear combination of $a_0,a_1,a_2$, and the coefficient of $s$ occurs only in the distinguished row.  Hence $T(s)$ has the form \eqref{eq:linearized-operator} for suitable matrices $C_i,D_i$, with the required row condition on the $D_i$.

Suppose now that $T(s)^N=0$ for some $N\ge1$.  In the matrix ring over the unital ring $\widetilde B[s]$ we have
\[
(1-T(s))Q(s)=1,
\qquad
Q(s)=1+T(s)+\cdots+T(s)^{N-1}.
\]
Consider the column of $Q(s)$ indexed by $0$.  Write
\[
q_0=Q_{0,0}(s),
\qquad
q_{\alpha,j}=Q_{(\alpha,j),0}(s).
\]
The $(\alpha,j),0$ entries of $(1-T)Q=1$ give, successively,
\begin{align*}
q_{\alpha,\ell_\alpha}
&=a_{i_{\alpha,\ell_\alpha}}q_0,\\
q_{\alpha,j}
&=a_{i_{\alpha,j}}q_{\alpha,j+1}
\qquad(2\le j<\ell_\alpha).
\end{align*}
Therefore
\begin{equation}\label{eq:qchain}
q_{\alpha,2}
=a_{i_{\alpha,2}}a_{i_{\alpha,3}}\cdots a_{i_{\alpha,\ell_\alpha}}q_0.
\end{equation}
The $0,0$ entry of $(1-T)Q=1$, together with \eqref{eq:T00}, \eqref{eq:Troot}, and \eqref{eq:qchain}, is
\[
q_0
=1+s\,f(a_0,a_1,a_2)q_0.
\]
The constant coefficient of this identity is $1$, so $q_0\ne0$.  Write
\[
q_0=b_0+b_1s+\cdots+b_es^e,
\qquad b_e\ne0.
\]
Since $s$ is central, comparison of coefficients gives
\[
b_0=1,
\qquad
b_{n+1}=f(a_0,a_1,a_2)b_n
\quad(0\le n<e),
\]
and the coefficient of $s^{e+1}$ gives
\[
f(a_0,a_1,a_2)b_e=0.
\]
Inductively $b_n=f(a_0,a_1,a_2)^n$, and hence
\[
f(a_0,a_1,a_2)^{e+1}=0.
\]
Thus $f(a_0,a_1,a_2)$ is nilpotent.
\end{proof}

For each nonzero $f\in\Free$, choose matrices $C_i(f),D_i(f)$ furnished by Lemma~\ref{lem:linearization}.  For a scalar triple $y=(y_0,y_1,y_2)\in k^3$ define
\begin{equation}\label{eq:scalar-pencil}
P_f(y;s)=\sum_{i=0}^2 y_i\bigl(C_i(f)+sD_i(f)\bigr)
\in\Mat_{d_f}(k[s]).
\end{equation}
The variable $s$ occurs only in the first row of this matrix.

\section{Degree bounds for products of the finite matrices}

The restriction on the $s$-coefficient gives the following degree bound.

\begin{lemma}[Degree and multihomogeneity of minors]\label{lem:minor-degree}
Let
\[
P(y;s)=\sum_{i=0}^2y_i(C_i+sD_i)\in\Mat_d(k[s]),
\]
where every $D_i$ is zero outside the first row.  Consider a product of $L$ matrices of this form.  Some factors may be evaluated at fixed vectors of $k^3$, and the remaining factors are evaluated at pairwise independent variable triples
\[
y^{(1)},\ldots,y^{(N)}.
\]
Assume that each variable triple occurs in exactly one factor.

Then every $r\times r$ minor $F(s)$ of the product satisfies:
\begin{enumerate}[label=\textup{(\roman*)}]
\item $\deg_sF\le L$ if $F\ne0$;
\item for every $h\ge0$, the coefficient $\coeff{h}{F}$ is multihomogeneous of degree $r$ in each variable triple $y^{(j)}$.
\end{enumerate}
\end{lemma}

\begin{proof}
For one factor, every matrix entry is linear in the three coordinates of $y$, so every $r\times r$ minor is homogeneous of degree $r$ in $y$.  Moreover $s$ occurs in at most one selected row.  In the determinant expansion, therefore, each term contains at most one entry involving $s$, and the degree in $s$ is at most one.

For a product, apply the Cauchy--Binet formula repeatedly.  Every $r\times r$ minor of a product is a sum of products of $r\times r$ minors, one from each factor.  Hence the $s$-degree of every summand is at most $L$.  If $y^{(j)}$ occurs in one factor, the minor taken from that factor is homogeneous of degree $r$ in $y^{(j)}$, while all other factors are independent of $y^{(j)}$.  Thus every summand is multihomogeneous of degree $r$ in each variable block.  Summation, and then taking a coefficient of $s$, preserves this property.  The zero polynomial causes no exception because it is multihomogeneous of every degree by convention.
\end{proof}

\section{Rank reduction with prescribed periodic values}\label{sec:rank}

Fix an integer $Q\ge1$, a subset
\[
S\subseteq\{0,1,\ldots,Q-1\},
\]
and a nonzero vector $c_r\in k^3$ for each $r\in S$.  At every integer position congruent to $r$ modulo $Q$, the value $c_r$ is prescribed.  No value is prescribed on the other residue classes.  A finite word
\[
w=(w_0,\ldots,w_{\ell-1})\in(k^3)^\ell
\]
is \emph{compatible with the prescribed values} if
\[
w_j=c_r
\quad\text{whenever }j\equiv r\pmod Q\text{ and }r\in S.
\]
Set
\[
H=Q-|S|.
\]
Thus $H$ is the number of residue classes modulo $Q$ on which no value has been prescribed.

The construction below assumes
\begin{equation}\label{eq:half-free}
|S|<\frac Q2,
\qquad\text{equivalently}\qquad Q<2H.
\end{equation}

\begin{lemma}[Rank factorization from a nonsingular minor]\label{lem:rank-factor}
Let $F$ be a field and let $A\in\Mat_d(F)$ have rank $r>0$.  Choose sets $I,J$ of $r$ row and column indices such that
\[
M=A_{I,J}
\]
is invertible.  Put
\[
U=A_{\bullet,J}\in\Mat_{d\times r}(F),
\qquad
V=A_{I,\bullet}\in\Mat_{r\times d}(F).
\]
Then
\begin{equation}\label{eq:rank-factor}
A=UM^{-1}V.
\end{equation}
Consequently, for every $B\in\Mat_d(F)$,
\begin{equation}\label{eq:sandwich}
ABA=UM^{-1}(VBU)M^{-1}V,
\end{equation}
and
\begin{equation}\label{eq:pivot-sandwich}
(ABA)_{I,J}=VBU.
\end{equation}
In particular, if $\det(ABA)_{I,J}=0$, then $\rank(ABA)<r$.
\end{lemma}

\begin{proof}
The $r$ columns of $U$ are linearly independent because their restriction to the rows $I$ is the invertible matrix $M$.  Since $\rank A=r$, these columns form a basis for the column space of $A$.  Therefore the $j$th column of $A$ is $U\lambda_j$ for a unique vector $\lambda_j\in F^r$.  Restricting to the rows $I$ gives
\[
V_{\bullet,j}=M\lambda_j,
\]
so $\lambda_j=M^{-1}V_{\bullet,j}$.  This proves \eqref{eq:rank-factor}.  Equation \eqref{eq:sandwich} follows by substitution.  Since $U_{I,\bullet}=M$ and $V_{\bullet,J}=M$, restricting \eqref{eq:sandwich} to rows $I$ and columns $J$ gives \eqref{eq:pivot-sandwich}.

If $\det(ABA)_{I,J}=0$, then \eqref{eq:pivot-sandwich} gives $\rank(VBU)<r$.  Equation \eqref{eq:sandwich} then implies
\[
\rank(ABA)\le\rank(VBU)<r.
\]
\end{proof}

For a finite word $w=(w_0,\ldots,w_{\ell-1})$, write
\[
P[w]=P(w_0;s)P(w_1;s)\cdots P(w_{\ell-1};s),
\]
with $P[\varnothing]=I_d$ for the empty word.  Concatenation of words is indicated by juxtaposition.

\begin{proposition}[A zero product with prescribed periodic values]\label{prop:zero-word}
Let $k$ be algebraically closed and let
\[
P(y;s)=\sum_{i=0}^2y_i(C_i+sD_i)\in\Mat_d(k[s]),
\]
where each $D_i$ is zero outside the first row.  Fix data as above, with values prescribed on fewer than half of the residue classes modulo $Q$.

Then there exists a nonempty finite word
\[
w=(w_0,\ldots,w_{\ell-1})
\]
such that
\begin{enumerate}[label=\textup{(\roman*)}]
\item every $w_j$ is nonzero;
\item $Q\mid\ell$;
\item $w$ is compatible with the prescribed values;
\item
\[
P(w_0;s)P(w_1;s)\cdots P(w_{\ell-1};s)=0
\quad\text{in }\Mat_d(k[s]).
\]
\end{enumerate}
\end{proposition}

\begin{proof}
Let $H=Q-|S|$, so $Q<2H$.  We construct words of length divisible by $Q$ that are compatible with all prescribed values while strictly lowering, over $k(s)$, the rank of their matrix products.

Start with the empty word $w$.  Its product is the identity matrix, so its rank is $d$.  Suppose at some stage that $w$ has length $\ell$, where $Q\mid\ell$, that it is compatible with all prescribed values and all its letters are nonzero, and that
\[
A=P[w]
\]
has positive rank
\[
r=\rank_{k(s)}A>0.
\]
Choose row and column sets $I,J$ of cardinality $r$ such that
\[
\det A_{I,J}\ne0.
\]

Choose an integer $m\ge1$ so large that
\begin{equation}\label{eq:m-choice}
2\ell+mQ+1\le2mH.
\end{equation}
Such an $m$ exists because this inequality is equivalent to
\[
2\ell+1\le m(2H-Q),
\]
and $2H-Q$ is a positive integer.

Consider a middle block of length $mQ$.  At every position whose residue modulo $Q$ is assigned, put the prescribed nonzero vector.  At every position whose residue class has no prescribed value, introduce a separate variable triple.  Since there are $H$ residue classes without prescribed values in each of the $m$ periods, the number of independent variable triples is exactly
\[
N=mH.
\]
Different positions belonging to the same unspecified residue class are treated as independent variable triples.

Let $c$ denote this variable middle block and consider the product corresponding to the concatenated word
\[
w\,c\,w.
\]
Define
\[
F(s)=\det\bigl(P[w\,c\,w]_{I,J}\bigr).
\]
By Lemma~\ref{lem:minor-degree},
\[
\deg_sF\le2\ell+mQ
\]
if $F\ne0$, and each coefficient of $F$ is multihomogeneous of degree $r$ in each of the $N=mH$ variable triples.  List all coefficients from degree $0$ through degree $2\ell+mQ$, inserting zero polynomials when necessary.  This gives
\[
q=2\ell+mQ+1
\]
multihomogeneous polynomials of degree $r$ in every block.  By \eqref{eq:m-choice},
\[
q\le2mH=2N.
\]
Lemma~\ref{lem:projective} therefore gives a specialization of every variable triple to a nonzero vector of $k^3$ for which all these coefficients vanish.  After this specialization, the middle block $c$ is an actual word of nonzero vectors, is compatible with all prescribed values, has length $mQ$, and satisfies
\begin{equation}\label{eq:pivot-killed}
\det\bigl(P[w\,c\,w]_{I,J}\bigr)=0
\quad\text{in }k[s].
\end{equation}

Put $B=P[c]$ and view $A,B$ as matrices over $k(s)$.  Since
\[
P[w\,c\,w]=ABA,
\]
Lemma~\ref{lem:rank-factor}, together with \eqref{eq:pivot-killed}, gives
\[
\rank_{k(s)}P[w\,c\,w]<r.
\]
Replace $w$ by $w\,c\,w$.  Its length is still divisible by $Q$.  Because the lengths of both the old word and the middle block are divisible by $Q$, concatenation does not change residue classes, so the new word remains compatible with all prescribed values.  Every letter is nonzero.

Thus every step strictly lowers a positive integer rank.  After at most $d$ steps the rank is zero.  A matrix over $k[s]$ of rank zero over $k(s)$ is the zero matrix, because $k[s]\hookrightarrow k(s)$ is injective.  The resulting word is nonempty and satisfies all four requirements.
\end{proof}

\section{A sequence with vanishing consecutive products}\label{sec:diagonal}

Let $k$ be a countable algebraically closed field.  The set of nonempty words in three letters is countable, so $\Free$ has a countable $k$-basis.  Since $k$ is countable, $\Free$ is countable.  Enumerate the nonzero elements of $\Free$ as
\[
f_1,f_2,f_3,\ldots,
\]
and let $P_j(y;s)=P_{f_j}(y;s)$ be the matrix defined in \eqref{eq:scalar-pencil}.

\begin{proposition}[A sequence with vanishing consecutive products]\label{prop:sequence}
There exists a sequence
\[
v_0,v_1,v_2,\ldots\in k^3\setminus\{0\}
\]
with the following property: for every $j$ there exists an integer $N_j\ge1$ such that for every $n\ge0$,
\begin{equation}\label{eq:uniform-window}
P_j(v_n;s)P_j(v_{n+1};s)\cdots P_j(v_{n+N_j-1};s)=0.
\end{equation}
\end{proposition}

\begin{proof}
Inductively prescribe nonzero vectors on residue classes modulo periods $Q_j$.  At every stage, fewer than half of the residue classes have prescribed values.

At stage $0$, take period $Q_0=1$ and assign nothing.

Suppose stage $j-1$ has period $Q_{j-1}$, with values prescribed on $a_{j-1}$ residue classes and
\[
\frac{a_{j-1}}{Q_{j-1}}<\frac12.
\]
Apply Proposition~\ref{prop:zero-word} to $P_j$ and the values prescribed at stage $j-1$.  We obtain a word compatible with those values
\[
w^{(j)}=(w^{(j)}_0,\ldots,w^{(j)}_{\ell_j-1})
\]
of nonzero vectors, with
\[
Q_{j-1}\mid\ell_j
\]
and
\begin{equation}\label{eq:wordzeroj}
P_j(w^{(j)}_0;s)\cdots P_j(w^{(j)}_{\ell_j-1};s)=0.
\end{equation}

Choose an integer $M_j\ge1$ so large that, with
\[
Q_j=M_jQ_{j-1},
\]
we have
\begin{equation}\label{eq:newperiodlarge}
Q_j\ge\ell_j
\end{equation}
and
\begin{equation}\label{eq:densitychoice}
\frac{a_{j-1}}{Q_{j-1}}+\frac{\ell_j}{Q_j}<\frac12.
\end{equation}
This is possible because the first term on the left is already strictly less than $1/2$.

Define the prescribed values modulo $Q_j$ as follows.  First copy every previously prescribed value to all of its lifts modulo $Q_j$: a residue class modulo $Q_j$ receives that value whenever its reduction modulo $Q_{j-1}$ was assigned.  Then prescribe
\[
w^{(j)}_0,\ldots,w^{(j)}_{\ell_j-1}
\]
at the first $\ell_j$ residue classes modulo $Q_j$.  There is no conflict because $w^{(j)}$ respects all values prescribed at stage $j-1$.

The lifted previously prescribed values occupy exactly $M_ja_{j-1}$ residue classes.  Adding the first $\ell_j$ classes creates at most $\ell_j$ additional assignments.  Hence the new number $a_j$ of assigned classes satisfies
\[
a_j\le M_ja_{j-1}+\ell_j.
\]
Dividing by $Q_j=M_jQ_{j-1}$ and using \eqref{eq:densitychoice} gives
\[
\frac{a_j}{Q_j}<\frac12.
\]
Thus the induction continues, and every value prescribed at one stage is retained at every later stage.

Define the full sequence $(v_n)$ as follows.  If the position $n$ is prescribed at some stage, let $v_n$ be that prescribed value; this is well defined because later stages preserve it.  If $n$ is never prescribed, set
\[
v_n=(1,0,0).
\]
Every $v_n$ is nonzero, and the full sequence agrees with every value prescribed at every stage.

Fix $j$.  At stage $j$, the word $w^{(j)}$ occupies the first $\ell_j$ positions of every period $Q_j$, and $\ell_j\le Q_j$.  Hence the final sequence contains $w^{(j)}$ beginning at every multiple of $Q_j$.

Let $n\ge0$, and let $b$ be the least multiple of $Q_j$ satisfying $b\ge n$.  Then
\[
0\le b-n<Q_j.
\]
Since $\ell_j\le Q_j$, the interval of indices
\[
b,b+1,\ldots,b+\ell_j-1
\]
is contained in
\[
n,n+1,\ldots,n+2Q_j-1.
\]
Therefore the matrix product over the latter interval contains, as a contiguous subproduct, the zero product \eqref{eq:wordzeroj}.  By associativity and absorption by zero,
\[
P_j(v_n;s)\cdots P_j(v_{n+2Q_j-1};s)=0.
\]
Thus \eqref{eq:uniform-window} holds with $N_j=2Q_j$.
\end{proof}

\section{Weighted backward shifts}\label{sec:shifts}

Let $(v_n)$ be the sequence from Proposition~\ref{prop:sequence}.  Put
\[
K=k(t),
\qquad
V=K^{\NN},
\]
the $K$-vector space of all sequences $u=(u(0),u(1),\ldots)$ with entries in $K$.  Multiplication in $\End_K(V)$ is composition: $(fg)(u)=f(g(u))$.
For $i=0,1,2$ define a $K$-linear endomorphism $a_i\in\End_K(V)$ by
\begin{equation}\label{eq:shift}
(a_i u)(n)=v_n(i)u(n+1).
\end{equation}
Let
\begin{equation}\label{eq:Adef}
A=k\langle a_0,a_1,a_2\rangle_+
\subseteq\End_K(V)
\end{equation}
be the nonunital $k$-subalgebra generated by these three endomorphisms.

We first connect products of the scalar matrices $P(v_n;s)$ with powers of the corresponding matrix of endomorphisms.

Let
\[
P(y;s)=\sum_{i=0}^2y_i(C_i+sD_i)\in\Mat_d(k[s])
\]
be any matrix of the form used above.  Define
\begin{equation}\label{eq:operator-pencil}
\widehat P(a;s)
=\sum_{i=0}^2(C_i+sD_i)a_i
\in\Mat_d(\End_K(V)[s]).
\end{equation}
Explicitly, its $(p,q)$ entry is
\[
\sum_{i=0}^2(C_i)_{pq}a_i
+s\sum_{i=0}^2(D_i)_{pq}a_i.
\]

\begin{lemma}[Coefficient action of a power]\label{lem:band}
For every $N\ge1$, every $h\ge0$, every matrix indices $p,q$, every $u\in V$, and every $n\ge0$,
\begin{equation}\label{eq:bandidentity}
\left(\coeff{h}{(\widehat P(a;s)^N)_{pq}}\,u\right)(n)
=
\left(\coeff{h}{\bigl(P(v_n;s)\cdots P(v_{n+N-1};s)\bigr)_{pq}}\right)u(n+N).
\end{equation}
On the right, the coefficient belongs to $k$ and is viewed inside $K$.
\end{lemma}

\begin{proof}
We argue by induction on $N$.

For $N=1$, the only possibly nonzero coefficients are those of $s^0$ and $s^1$.  By \eqref{eq:shift},
\begin{align*}
\left(\sum_i(C_i)_{pq}a_i\right)u(n)
&=\left(\sum_i v_n(i)(C_i)_{pq}\right)u(n+1),\\
\left(\sum_i(D_i)_{pq}a_i\right)u(n)
&=\left(\sum_i v_n(i)(D_i)_{pq}\right)u(n+1),
\end{align*}
which is precisely \eqref{eq:bandidentity}.

Assume the formula holds for $N$.  Since
\[
\widehat P^{N+1}=\widehat P\,\widehat P^N,
\]
the coefficient of $s^h$ in its $(p,q)$ entry is
\[
\sum_{r=1}^d\sum_{a+b=h}
\coeff{a}{\widehat P_{pr}}\,
\coeff{b}{(\widehat P^N)_{rq}},
\]
where multiplication in $\End_K(V)$ is composition.  Apply this endomorphism to $u$ and evaluate at $n$.  Since multiplication in $\End_K(V)$ is composition, the right factor acts on $u$ first; the left factor then evaluates that result at $n+1$.  Using the case $N=1$ for the left factor and the induction hypothesis for the right factor gives
\[
\left(\coeff{a}{P(v_n;s)_{pr}}\right)
\left(\coeff{b}{(P(v_{n+1};s)\cdots P(v_{n+N};s))_{rq}}\right)u(n+N+1).
\]
Summing over the intermediate matrix index $r$ and over $a+b=h$ is exactly the coefficient rule for the matrix product
\[
P(v_n;s)P(v_{n+1};s)\cdots P(v_{n+N};s).
\]
This proves \eqref{eq:bandidentity} for $N+1$.
\end{proof}

\begin{corollary}\label{cor:window-nilpotent}
If for some $N\ge1$
\[
P(v_n;s)\cdots P(v_{n+N-1};s)=0
\qquad\text{for every }n\ge0,
\]
then
\[
\widehat P(a;s)^N=0
\quad\text{in }\Mat_d(\End_K(V)[s]).
\]
\end{corollary}

\begin{proof}
Every coefficient of every entry of $\widehat P(a;s)^N$ acts as zero on every $u\in V$ at every coordinate $n$ by Lemma~\ref{lem:band}.  Hence each coefficient is the zero endomorphism, so the polynomial matrix itself is zero.
\end{proof}

\section{The algebra generated by the shifts is nil}\label{sec:nil}

\begin{proposition}\label{prop:A-nil}
The algebra $A$ defined in \eqref{eq:Adef} is nil.
\end{proposition}

\begin{proof}
Let $x\in A$.  If $x=0$, there is nothing to prove.  Assume $x\ne0$.  By definition of the nonunital algebra generated by $a_0,a_1,a_2$, there exists a nonzero $f\in\Free$ such that
\[
x=f(a_0,a_1,a_2).
\]
Choose $j$ with $f_j=f$ in the enumeration of Section~\ref{sec:diagonal}.  Let $P_j=P_f$ be its scalar linearization matrix.

By Proposition~\ref{prop:sequence}, there is $N_j\ge1$ such that every consecutive product of length $N_j$ of $P_j(v_n;s)$ is zero.  Corollary~\ref{cor:window-nilpotent} therefore shows that the corresponding matrix of endomorphisms
\[
T_f(a;s)=\widehat P_j(a;s)
\]
is nilpotent.  Lemma~\ref{lem:linearization} now implies that
\[
f(a_0,a_1,a_2)=x
\]
is nilpotent.  Since $x$ was arbitrary, $A$ is nil.
\end{proof}

\section{A nonnilpotent two-by-two matrix over the nil algebra}\label{sec:matrix}

For each $n\ge0$, define
\begin{equation}\label{eq:pn}
p_n(t)=v_n(0)+t\,v_n(1)+t^2v_n(2)\in k[t].
\end{equation}
Because $v_n\ne0$, at least one of its three coordinates is nonzero, so $p_n(t)$ is a nonzero polynomial.  Hence it is invertible in $K=k(t)$.

Define $z\in V$ by
\begin{equation}\label{eq:z}
z(n)=\left(\prod_{j=0}^{n-1}p_j(t)\right)^{-1},
\end{equation}
with the empty product equal to $1$.  In particular $z(0)=1$, so $z\ne0$.  Moreover
\[
p_n(t)z(n+1)=z(n),
\]
and therefore, by \eqref{eq:shift},
\begin{equation}\label{eq:fixed}
(a_0+t a_1+t^2a_2)z=z.
\end{equation}

Since $A$ is nil and $a_0\in A$, choose $m\ge1$ with
\[
a_0^m=0.
\]
Inside the unital ring $\End_K(V)$ put
\[
g=1+a_0+a_0^2+\cdots+a_0^{m-1}.
\]
Then
\begin{equation}\label{eq:geom}
g(1-a_0)=1-a_0^m=1.
\end{equation}
Define
\begin{equation}\label{eq:bc}
b=ga_1,
\qquad
c=ga_2.
\end{equation}
Although $g$ need not belong to the nonunital algebra $A$, both $b$ and $c$ do, because
\[
b=a_1+a_0a_1+\cdots+a_0^{m-1}a_1,
\qquad
c=a_2+a_0a_2+\cdots+a_0^{m-1}a_2.
\]

Rearranging \eqref{eq:fixed} gives
\[
(1-a_0)z=t a_1z+t^2a_2z.
\]
Apply $g$ and use \eqref{eq:geom}.  Since $g$ is $K$-linear,
\[
z=t\,bz+t^2cz.
\]
Thus
\begin{equation}\label{eq:companion-relation}
bz+tcz=t^{-1}z.
\end{equation}

Consider the $K$-linear endomorphism of $V\oplus V$ represented by
\[
H=
\begin{pmatrix}
b&c\\
1&0
\end{pmatrix}.
\]
This matrix need not belong to $\Mat_2(A)$, but its square will.  Put
\[
\xi=
\begin{pmatrix}
z\\tz
\end{pmatrix}.
\]
Then \eqref{eq:companion-relation} gives
\[
H\xi
=
\begin{pmatrix}
bz+tcz\\z
\end{pmatrix}
=
\begin{pmatrix}
t^{-1}z\\z
\end{pmatrix}
=t^{-1}\xi.
\]
Hence
\[
H^2\xi=t^{-2}\xi.
\]
But
\begin{equation}\label{eq:W}
W:=H^2
=
\begin{pmatrix}
b^2+c&bc\\
b&c
\end{pmatrix}.
\end{equation}
Every entry of $W$ belongs to $A$, so
\[
W\in\Mat_2(A).
\]
Since $\xi\ne0$ and $t^{-2}\ne0$, for every $r\ge1$,
\[
W^r\xi=t^{-2r}\xi\ne0.
\]
Therefore $W$ is not nilpotent.  Combined with Proposition~\ref{prop:A-nil}, this proves the first assertion of Theorem~\ref{thm:main}.

\section{Two nil left ideals whose sum is not nil}\label{sec:kothe}

Set
\[
R=k\,1+A\subseteq\End_K(V).
\]
The identity endomorphism belongs to $R$, and closure under multiplication follows from
\[
(\lambda1+a)(\mu1+b)=\lambda\mu1+(\lambda b+\mu a+ab).
\]
Thus $R$ is a unital $k$-algebra.  Moreover $A$ is a two-sided ideal of $R$: if $r=\lambda1+a\in R$ and $x\in A$, then
\[
rx=\lambda x+ax\in A,
\qquad
xr=\lambda x+xa\in A.
\]
The ideal $A$ is nil by Proposition~\ref{prop:A-nil}.

Let
\[
S=\Mat_2(R).
\]
Define
\[
L_1=
\left\{
\begin{pmatrix}
x&0\\y&0
\end{pmatrix}:x,y\in A
\right\},
\qquad
L_2=
\left\{
\begin{pmatrix}
0&x\\0&y
\end{pmatrix}:x,y\in A
\right\}.
\]
For
\[
M=\begin{pmatrix}r_{11}&r_{12}\\r_{21}&r_{22}\end{pmatrix}\in S
\]
and $x,y\in A$, one has
\[
M\begin{pmatrix}x&0\\y&0\end{pmatrix}
=
\begin{pmatrix}r_{11}x+r_{12}y&0\\r_{21}x+r_{22}y&0\end{pmatrix}\in L_1
\]
and
\[
M\begin{pmatrix}0&x\\0&y\end{pmatrix}
=
\begin{pmatrix}0&r_{11}x+r_{12}y\\0&r_{21}x+r_{22}y\end{pmatrix}\in L_2,
\]
because $A$ is a left ideal of $R$.  Thus $L_1$ and $L_2$ are left ideals of $S$.

They are nil.  Indeed, if
\[
X=
\begin{pmatrix}
x&0\\y&0
\end{pmatrix}\in L_1,
\]
then induction gives
\[
X^n=
\begin{pmatrix}
x^n&0\\yx^{n-1}&0
\end{pmatrix}
\qquad(n\ge1).
\]
If $x^m=0$, then $X^{m+1}=0$.  Similarly, for
\[
Y=
\begin{pmatrix}
0&x\\0&y
\end{pmatrix}\in L_2,
\]
one has
\[
Y^n=
\begin{pmatrix}
0&xy^{n-1}\\0&y^n
\end{pmatrix}
\qquad(n\ge1),
\]
so $y^m=0$ implies $Y^{m+1}=0$.

Finally,
\[
L_1+L_2=\Mat_2(A),
\]
and this sum contains the nonnilpotent matrix $W$ from \eqref{eq:W}.  Thus $L_1$ and $L_2$ are nil left ideals of the unital ring $S$, while $L_1+L_2$ is not nil.  This is a direct counterexample to K\"othe's conjecture and completes the proof of Theorem~\ref{thm:main}.

\begin{thebibliography}{99}

\bibitem{Adamczewski2026}
T.~Adamczewski (maintainer),
\emph{K\"othe conjecture: disproof (Krempa's matrix form)},
Lean~4 formalization repository, 2026, commit \texttt{b05275587973}.
\url{https://github.com/tadamcz/koethe/commit/b0527558797325e98a8e57f2197e527210290dc4}.

\bibitem{AmitsurInfinite1956}
S.~A.~Amitsur,
Algebras over infinite fields,
\emph{Proc. Amer. Math. Soc.} \textbf{7} (1956), 35--48.
\url{https://doi.org/10.1090/S0002-9939-1956-0075933-2}.

\bibitem{AmitsurRadicals1956}
S.~A.~Amitsur,
Radicals of polynomial rings,
\emph{Canad. J. Math.} \textbf{8} (1956), 355--361.
\url{https://doi.org/10.4153/CJM-1956-040-9}.

\bibitem{BeidarFongPuczylowski2001}
K.~I.~Beidar, Y.~Fong, and E.~R.~Puczy\l{}owski,
Polynomial rings over nil rings cannot be homomorphically mapped onto rings with nonzero idempotents,
\emph{J. Algebra} \textbf{238} (2001), no.~1, 389--399.
\url{https://doi.org/10.1006/jabr.2000.8632}.

\bibitem{Hartshorne}
R.~Hartshorne,
\emph{Algebraic Geometry},
Graduate Texts in Mathematics, vol.~52,
Springer-Verlag, New York, 1977.

\bibitem{Kothe1930}
G.~K\"othe,
Die Struktur der Ringe, deren Restklassenring nach dem Radikal vollst\"andig reduzibel ist,
\emph{Math. Z.} \textbf{32} (1930), 161--186.
\url{https://doi.org/10.1007/BF01194626}.

\bibitem{Krempa1972}
J.~Krempa,
Logical connections between some open problems concerning nil rings,
\emph{Fund. Math.} \textbf{76} (1972), no.~2, 121--130.
\url{https://doi.org/10.4064/fm-76-2-121-130}.

\bibitem{Sands1973}
A.~D.~Sands,
Radicals and Morita contexts,
\emph{J. Algebra} \textbf{24} (1973), no.~2, 335--345.
\url{https://doi.org/10.1016/0021-8693(73)90143-9}.

\bibitem{Smoktunowicz2000}
A.~Smoktunowicz,
Polynomial rings over nil rings need not be nil,
\emph{J. Algebra} \textbf{233} (2000), no.~2, 427--436.
\url{https://doi.org/10.1006/jabr.2000.8451}.

\bibitem{Smoktunowicz2005}
A.~Smoktunowicz,
On primitive ideals in polynomial rings over nil rings,
\emph{Algebr. Represent. Theory} \textbf{8} (2005), 69--73.
\url{https://doi.org/10.1007/s10468-004-6118-7}.

\end{thebibliography}

\end{document}
